\icode{E-MK-MU-L3-260625-002}
物理学中加速度的定义式为\(a = \frac{v_t - v_0}{t}\)，而历史上有些科学家曾把相等位移内速度变化相等的单向直线运动称为“匀变速直线运动”（现称“另类匀变速直线运动”），“另类加速度”定义为\(A = \frac{v_t - v_0}{x}\)，其中\(v_t\)和\(v_0\)分别表示某段位移\(x\)内的初速度和末速度。\(A > 0\)表示物体做加速运动，\(A < 0\)表示物体做减速运动。则下列说法正确的是\paren[CD]
\begin{choices}
  \item 若\(A\)不变，则\(a\)也不变
  \item 若\(A\)不变，则物体在位移中点处的速度比\(\frac{v_0 + v_t}{2}\)大
  \item 若\(a\)不变，则物体在中间时刻的速度为\(\frac{v_0 + v_t}{2}\)
  \item 若\(a > 0\)且保持不变，则\(A\)逐渐变小
\end{choices}
\begin{answer}
  CD
\end{answer}